% !TEX encoding = UTF-8
% !TEX TS-program = pdflatex
% !TEX root = ../tesi.tex

%**************************************************************
\chapter{Introduzione}
\label{cap:introduzione}
%**************************************************************

%Introduzione al contesto applicativo.\\

%\noindent Esempio di utilizzo di un termine nel glossario \\
%\gls{api}. \\

%\noindent Esempio di citazione in linea \\
%\cite{site:agile-manifesto}. \\

%\noindent Esempio di citazione nel pie' di pagina \\


%**************************************************************
\section{L'azienda: SyncLab s.r.l.}


\begin{figure}[!h] 
	\centering 
	\includegraphics[width=0.25\columnwidth]{immagini/logo_azienda} 
	\caption{Logo dell'azienda ospitante}
	\label{figura:logo-azienda}
\end{figure} 
SyncLab nasce nel 2002 come Software house, trasformandosi poi nella figura di \gls{System Integrator} in conseguenza del rapido ampliamento delle competenze e del dominio tecnologico offerto; oggi SyncLab vanta un organico di oltre 200 persone e 4 sedi sul territorio italiano (Napoli, Roma, Padova e Milano).\\
La grande attenzione alla gestione delle risorse umane ha fatto di Sync Lab un
riferimento in positivo per quanti volessero avviare o far evolvere in chiave professionale
la propria carriera.
Il basso turn-over testimonia la voglia dei collaboratori di condividere il progetto
comune, assumendo all’interno di esso ruoli e responsabilità che solo un processo
evolutivo così intenso può offrire.
I ricavi hanno avuto un incremento proporzionale alla crescita dell’azienda beneficiando
dell’approccio adattivo e diversificato al mercato


%**************************************************************
\subsection{Serivzi offerti}
Sync Lab Srl è un’azienda leader nella consulenza tecnologica, impegnata in un processo
continuo di identificazione e messa in opera di soluzioni per i clienti finalizzate alla
creazione di valore. Essa supporta le esigenze di innovazione di tutte le organizzazioni
ed in ogni settore di mercato nell’ambito Information Technology, con servizi in ambito:
\begin{itemize}
	\item \textit{Business Consultancy};
	\item \textit{Project Financing};
	\item \textit{IT Consultancy}.
\end{itemize}
L’azienda ha come punti di forza la qualità dei servizi offerti (certificazioni \gls{ISO} 9001,
ISO 14001, ISO 27001, OHSAS 18001) ed un’accurata gestione delle risorse umane.
L’approfondita conoscenza di processi e tecnologie, maturata in esperienze altamente
significative e qualificanti le fornisce l’expertise e Know How necessari per gestire
progetti di elevata complessità, dominando l’intero ciclo di vita: Studio di fattibilità,
Progettazione, Implementazione, Governance e Post Delivery.
L’offerta di consulenza specialistica trova le punte di eccellenza nella progettazione
di architetture Software avanzate, siano esse per applicativi di dominio, per sistemi
di supporto al business, per sistemi di integrazione o per sistemi di monitoraggio
applicativo/territoriale.
Il suo laboratorio RD è sempre al passo con i nuovi paradigmi tecnologici e di comu-
nicazione (\textit{\gls{Big Data Analysis}}, \textit{\gls{Cloud Computing}}, \textit{\gls{Internet delle Cose}},\textit{ Mobile}
e Sicurezza IT) per supportare i propri clienti nella creazione ed integrazione di
applicazioni, processi e dispositivi.
Le attività in ambito Educational ed RD hanno permesso di acquisire una profonda
conoscenza degli strumenti di finanza agevolata fruendone direttamente ed interagendo
con enti di supporto ai progetti innovativi dei propri clienti. L’azienda, grazie alla rete
di relazioni a livello nazionale ed internazionale, ha ottenuto importanti finanziamenti
in progetti RD (FP7 e H2020) della Comunità Europea.
Sync Lab si sta sempre più specializzando in vari settori d’impiego: dal settore bancario
a quello assicurativo con una nicchia importante nell’ambito sanità, in cui vanta
un prodotto d’eccellenza per la gestione delle cliniche private. L’azienda inoltre ha
recentemente fondato una collegata Sync Security che si occupa espressamente del
mondo della cyber security e sicurezza informatica in genere.
%**************************************************************
